\chapter{Visual Sparsity via Decoder-Encoder Cross-Attention: Why Attention Rollout Does Not Apply (Day 6)}
\label{chap:day06}

\section{Motivation}

Days 3--5 characterized the model's grounded boxes and tested whether masking the top-ranked one
changes the answer, but neither day asked the more direct visual-attribution question: when
Florence-2 grounds a phrase, where in the image does it actually attend while doing so? Visual
sparsity is one of the six metrics defined for evaluating black-box explainable-AI frameworks
\citep{arreche2024exai}, and it answers exactly that question --- given a per-pixel or per-patch
attribution signal, how concentrated is the attention mass that signal assigns across the image?
A sparse, concentrated heatmap is the visual signature of a focused explanation; a heatmap spread
thinly and evenly across the frame carries comparatively little information about where the
judgment actually came from. Day 6's task was to produce that attribution signal for Florence-2
and score its concentration across the pilot's samples. The plan, written before this day's
implementation began, specified a particular method for producing the signal: attention rollout.
That plan did not survive contact with Florence-2's actual architecture, and the deviation that
resulted --- discovered, verified, and resolved before any code was written against the wrong
assumption --- is this chapter's centerpiece.

\section{Method Deviation from the Plan}
\label{sec:day06-deviation}

The plan called for \texttt{attention\_rollout()} first, the textbook ViT-style explainability
method that recursively composes self-attention matrices across the layers of a single-modality
self-attention-only encoder to estimate how much each input token ultimately propagates into each
output token \citep{abnar2020rollout}. Before implementing it, Florence-2's actual cached modeling
code (\texttt{modeling\_florence2.py}) was checked directly, rather than assuming a textbook
ViT-style rollout would apply unmodified --- and it does not apply cleanly.

\textbf{The verified architecture fact.} Florence-2's encoder does not run a self-attention stack
over image patches alone. It runs \textbf{one shared self-attention stack over a fused
image+text sequence}: \texttt{[1 global ``spatial\_avg\_pool'' token] + [576 image-patch tokens,
24$\times$24] + [text prompt tokens]}. This was confirmed directly against the model's own code,
not inferred from documentation or assumed by analogy to a generic vision transformer. Two
specific checks established it:

\begin{itemize}
\item \texttt{\_encode\_image} returns \texttt{torch.cat([spatial\_avg\_pool\_x,
temporal\_avg\_pool\_x])}, which is $1 + 576 = 577$ tokens, confirmed by printing
\texttt{image\_features.shape} directly rather than reading the count off a docstring.
\item $576 = 24 \times 24$ follows from the DaViT vision tower's four convolutional stages
(strides 4, 2, 2, 2) applied to the processor's fixed $768 \times 768$ input image: $768 \to 192
\to 96 \to 48 \to 24$, computed by hand from the stride arithmetic rather than taken on faith.
\end{itemize}

\textbf{Why this breaks rollout.} Classic attention rollout \citep{abnar2020rollout} is defined
for a self-attention-only stack processing \emph{one} modality: the method's recursive
matrix-composition step assumes that what is being mixed across layers is attention among tokens
of a single, homogeneous kind, so that the result can be read as ``how much does input patch $i$
ultimately influence output patch $j$.'' Florence-2's encoder violates that assumption at its
foundation: every self-attention layer in the fused stack mixes image-patch tokens with text
prompt tokens at the same time, in the same attention matrix. Rolling out Florence-2's encoder
as if it were a single-modality stack would recursively mix in attention to text tokens at every
layer in a way the method was never designed for, and that mixing is not separable back out after
the fact --- there is no post-hoc operation that can take the rolled-out, jointly-mixed result and
cleanly attribute the image-only portion of it. The architecture was checked first, specifically
because the plan's chosen method carries this hidden precondition, and the check showed the
precondition fails.

\textbf{What was used instead.} \textbf{Decoder$\to$encoder cross-attention} was substituted for
rollout. For each token the decoder generates, it attends back to the full 577-token encoder
sequence; \texttt{attribution.py}'s \texttt{cross\_attention\_heatmap()} averages this
cross-attention over decoder layers, heads, and generated tokens, then drops the global token
(index 0) before reshaping the remaining 576 values into the 24$\times$24 patch grid. Unlike
rollout, cross-attention requires no recursive composition and no single-modality assumption: it
is a per-step, already-computed quantity --- ``where did the decoder look, this step, to produce
this output token'' --- read directly off the model rather than approximated by composing layers.
This substitution is well-established for encoder-decoder transformers generally (cross-attention
visualization has long been used for machine-translation attribution) and sidesteps the fused-modality
problem entirely, because cross-attention is always computed against the same fixed encoder
sequence regardless of how that sequence's internal self-attention mixed modalities together.

Other attribution families were also considered and set aside for this architecture, a point
returned to in Section~\ref{sec:day06-implications}: gradient- and perturbation-based attribution
methods of the kind implemented in Captum \citep{kokhlikyan2020captum}, layer-wise relevance
propagation \citep{bach2015lrp}, and DeepLIFT \citep{shrikumar2017deeplift} all could, in
principle, be adapted to Florence-2, but each carries its own architectural assumptions that would
need the same kind of verification rollout was given here, and none was implemented or verified
within this pilot's scope.

\section{Architecture: An Attribution Layer}

Figure~\ref{fig:diagram-d} lays out the attribution pipeline that resulted from
Section~\ref{sec:day06-deviation}'s deviation. Generation for attribution purposes always runs in
greedy mode (\texttt{num\_beams=1}), a difference from Days 3--5 taken up in
Section~\ref{sec:day06-implementation}. The decoder's cross-attention to the full 577-token
encoder sequence is captured at every generation step, averaged over layers, heads, and generated
steps into a single 577-length vector, then the global token is dropped and the remaining 576
values are reshaped into the 24$\times$24 patch grid and normalized to $[0,1]$ before the sparsity
metrics are computed on it.

\begin{figure}[h]
\centering
\resizebox{\textwidth}{!}{%
\begin{tikzpicture}[node distance=1.5cm]
\node[block] (gen) {Image + prompt\\$\to$ decoder token $t$\\(greedy, num\_beams=1)};
\node[block, right=of gen] (cross) {Cross-attention\\to 577-token\\encoder sequence};
\node[block, right=of cross] (avg) {Average over\\layers, heads,\\generated steps};
\node[block, below=of avg] (drop) {Drop global token,\\reshape 576 $\to$ 24$\times$24};
\node[block, left=of drop] (norm) {Normalize\\to $[0,1]$};
\node[block, left=of norm] (metric) {Sparsity metrics:\\topk\_mass\_ratio,\\regions\_above\_threshold};
\draw[arrow] (gen) -- (cross);
\draw[arrow] (cross) -- (avg);
\draw[arrow] (avg) -- (drop);
\draw[arrow] (drop) -- (norm);
\draw[arrow] (norm) -- (metric);
\end{tikzpicture}%
}
\caption{Attribution layer: decoder$\to$encoder cross-attention, substituted for attention rollout.}
\label{fig:diagram-d}
\end{figure}

Nothing about this pipeline reaches back into Florence-2's encoder self-attention --- the
fused-modality problem identified in Section~\ref{sec:day06-deviation} is avoided entirely by
operating only on the cross-attention edge between decoder and encoder, a quantity the fused
encoder's internal mixing does not touch.

\section{Implementation Details}
\label{sec:day06-implementation}

\texttt{attribution.py}'s \texttt{cross\_attention\_heatmap()} implements
Figure~\ref{fig:diagram-d}'s pipeline directly. Two implementation details were significant enough
to flag explicitly rather than leave implicit in the code.

\textbf{The off-by-one trap.} 577 is not a perfect square. Computing
$\sqrt{\texttt{num\_image\_tokens}}$ directly --- the natural first move for reshaping a sequence
of attention weights into a square grid --- would have silently misaligned every heatmap by one
cell, because the result, $\sqrt{577} \approx 24.02$, is not an integer and gives no warning that
something is wrong before producing a plausible-looking but wrong grid. The actual fix relies on
the architecture fact verified in Section~\ref{sec:day06-deviation}: index 0 of the 577-token
sequence is the global \texttt{spatial\_avg\_pool} token, not a patch, so it must be dropped
\emph{before} reshaping, leaving exactly 576 values for a clean $24 \times 24$ grid. This is
flagged here as a trap rather than a routine implementation step because it would have failed
silently --- the reshape would not raise an error on a near-square count, it would simply produce
a heatmap shifted by one cell relative to the true patch layout, with no symptom visible short of
comparing the result against an independently verified ground truth.

\textbf{Decoding-mode difference from Chapters 3--5.} Cross-attention extraction always uses
\texttt{num\_beams=1} (greedy decoding), not the \texttt{num\_beams=3} beam search used for the
logged answers in Chapters 3--5. This is a deliberate substitution, not an oversight: Hugging
Face's beam search reorders the batch dimension at every generation step to keep the surviving
beams contiguous, and \texttt{generate()}'s returned attentions are not automatically
un-reordered to compensate, so beam-search cross-attentions do not line up sample-by-sample with
the tokens they nominally belong to without extra bookkeeping this pilot's scope did not require.
Greedy decoding sidesteps the reordering problem entirely, at the cost of attributing a freshly
greedy-decoded output rather than the original beam-search answer. The substitution was
spot-checked rather than assumed safe: on image \texttt{0000007}, the greedy-decoded boxes
matched Chapter 4's beam-search boxes closely in position. The substitution changes \emph{how}
the output was decoded, not \emph{what} gets attributed --- the same regions are grounded either
way, which is what makes the cross-attention heatmap a meaningful attribution signal for the
beam-search answers reported elsewhere in this pilot, even though it is computed against a
greedy rerun.

\section{Results}

\texttt{scripts/06\_visual\_sparsity.py} ran the attribution pipeline over all 163 pilot samples.
4 of the 163 samples are Day 4's grid-fallback cases --- no model-returned box at all, and
therefore no grounded phrase to attribute attention to --- and were excluded from scoring, the
same exclusion Day 4 itself flagged. Sparsity was scored on the remaining \textbf{159 samples}.

\begin{center}
\begin{tabular}{lr}
\toprule
\textbf{Metric} & \textbf{Mean (159 scored samples)} \\
\midrule
Top-5-of-576-cells mass ratio & \textbf{0.108} \\
Top-20-of-576-cells mass ratio & 0.247 \\
Cells above 0.5 threshold (of 576) & 7.0 \\
\bottomrule
\end{tabular}
\end{center}

A top-5 mass ratio of 0.108 means the five highest-attention cells, out of 576 total, together
carry only 10.8\% of the total attention mass on average --- a long way from the concentration a
genuinely sparse, focused explanation would show, and a useful baseline against which the
by-rule and by-phrase breakdowns below can be read.

\subsection{By assigned rule}

\begin{center}
\begin{tabular}{lr}
\toprule
\textbf{Assigned rule} & \textbf{Top-5 mass ratio} \\
\midrule
rule\_1 (hard hat) & 0.125 \\
rule\_2 (harness) & 0.113 \\
rule\_3 (guardrail) & 0.098 \\
rule\_4 (struck-by risk) & 0.079 \\
\bottomrule
\end{tabular}
\end{center}

\subsection{By attributed phrase}

\begin{center}
\begin{tabular}{lr}
\toprule
\textbf{Attributed phrase} & \textbf{Top-5 mass ratio} \\
\midrule
hard hat & \textbf{0.195} \\
safety harness & 0.118 \\
worker & 0.115 \\
guardrail & 0.099 \\
excavator & \textbf{0.079} \\
\bottomrule
\end{tabular}
\end{center}

Hard hat is the most concentrated phrase in the breakdown, at 0.195; excavator is the least
concentrated, at 0.079 --- the two phrases anchor opposite ends of the by-phrase ordering, and the
gap between them is the subject of this chapter's central finding.

\subsection{Finding: concentration is driven mainly by target physical size, not explanation quality}

Hard hat's top-5 mass ratio, 0.195, is roughly \textbf{2.5$\times$} excavator's, 0.079. The
natural first reading of that gap is that the model's explanations for hard-hat judgments are
substantially better-focused than its explanations for excavator-related judgments. The data does
not support that reading once the grounded objects' own physical sizes are brought in. Mean
detected box area per phrase, taken from \texttt{baseline\_predictions.csv}, tracks the
concentration ordering closely:

\begin{center}
\begin{tabular}{lrr}
\toprule
\textbf{Phrase} & \textbf{Mean box area (px\textsuperscript{2})} & \textbf{Top-5 mass ratio} \\
\midrule
excavator & 287{,}913 & 0.079 \\
guardrail & 310{,}291 & 0.099 \\
safety harness & 178{,}491 & 0.118 \\
worker & 79{,}005 & 0.115 \\
hard hat & 94{,}330 & 0.195 \\
\bottomrule
\end{tabular}
\end{center}

The two largest objects, guardrail and excavator, are also the two least concentrated; the
smallest object, hard hat, is the most concentrated. The correlation between log(box area) and
top-5 mass ratio, computed across all scored samples, is \textbf{$-0.70$} --- a strong,
negative relationship between how physically large the grounded object is and how concentrated
the attention attributed to it appears.

This is a real confound, not a bug, and not a flaw in the attribution pipeline built in this
chapter. A small object's 576-cell patch grid mechanically has fewer cells it \emph{could} cover
in the first place: a hard hat that occupies a handful of patches will, almost by construction,
show its attention mass concentrated into a handful of cells, regardless of whether the
underlying reasoning that produced that attention is any good. A guardrail spanning much of the
frame has many more cells available to it geometrically, so even attention that is, in some
qualitative sense, just as well-targeted will register as more spread out by a metric defined
purely over cell counts. Visual sparsity, as measured here, partly reflects target footprint size
rather than purely reflecting explanation focus, and that confound has to be read into any later
comparison of sparsity across rules that attribute to objects of different physical sizes.

\section{Worked Examples}

Figure~\ref{fig:day06-worked} shows the two saved overlays that anchor the by-phrase finding
visually, one from each end of the concentration ordering.

\begin{figure}[h]
\centering
\begin{subcaptionblock}{0.48\linewidth}
\centering
\includegraphics[width=\linewidth]{sparsity/0000023_rule_1_hard_hat.png}
\caption{Image 23, rule~1 (hard hat): a small, sharp hotspot on a distant worker's hard hat.}
\end{subcaptionblock}
\hfill
\begin{subcaptionblock}{0.48\linewidth}
\centering
\includegraphics[width=\linewidth]{sparsity/0000079_rule_3_guardrail.png}
\caption{Image 79, rule~3 (guardrail): attention spread broadly across most of the frame.}
\end{subcaptionblock}
\caption{The clearest focused and scattered examples in the sample, consistent with the
hard-hat/excavator concentration gap and its physical-size explanation.}
\label{fig:day06-worked}
\end{figure}

Image 23's heatmap is the clearest ``focused'' example in the sample: a small, sharp hotspot
lands exactly on a distant worker's hard hat, against an otherwise empty dirt lot with nothing
nearby to compete for attention mass. Image 79's heatmap is the clearest ``scattered'' example:
attention spreads broadly across most of the frame, consistent with the guardrail's large,
elongated structural footprint rather than with any failure of the model to localize it. Read
side by side, the two panels make the size confound visually concrete --- the hard hat is small
and the heatmap is concentrated; the guardrail is large and the heatmap is not, and nothing about
either heatmap looks like a worse or better explanation in any qualitative sense, only a smaller
or larger target.

\section{Standard vs. Non-Standard Choices}

\textbf{Standard}: using decoder$\to$encoder cross-attention as an attribution signal for an
encoder-decoder transformer is itself well-established practice, the same basic technique behind
attention visualization for machine translation --- read the existing cross-attention weights
directly off the model's forward pass, average appropriately, and treat them as evidence of where
the model was looking when it produced a given output. Adopting that established technique here,
rather than inventing a new mechanism for reading attention out of Florence-2, follows precedent.

\textbf{Non-standard}: the explicit, verified decision \emph{not} to use attention rollout --- the
textbook ViT XAI method the plan originally specified --- is the non-standard, and reportable,
choice in this chapter. Most pipelines that adopt a planned method do so without re-checking
whether the target model's actual architecture satisfies that method's assumptions; this pilot
checked Florence-2's cached modeling code first, found that its fused image+text encoder breaks
attention rollout's single-modality precondition, and substituted cross-attention instead, with
the substitution's correctness established by direct inspection of the model's code rather than
by analogy or convention. That the substitution was discovered, documented, and resolved before
any rollout code was written against the wrong assumption is the genuine methodological
deviation this chapter records.

\section{Implications for the Paper}
\label{sec:day06-implications}

The architecture verification in Section~\ref{sec:day06-deviation} should be written up as a
candidate dedicated subsection in the paper on \textbf{adapting attribution methods to
fused-modality encoders}, not folded silently into an implementation appendix. The general lesson
generalizes well beyond Florence-2: any vision-language model whose encoder processes image and
text tokens through a single shared self-attention stack, rather than through separate
modality-specific towers, will break attention rollout's single-modality assumption in the same
way, for the same reason, and the fix --- falling back to a decoder$\to$encoder cross-attention
signal that does not require composing self-attention layers at all --- is also general, wherever
the architecture provides a decoder that attends back to the encoder's output.

This generalization extends past rollout to the other attribution families considered and set
aside in Section~\ref{sec:day06-deviation}. Captum's perturbation- and gradient-based methods
\citep{kokhlikyan2020captum}, layer-wise relevance propagation \citep{bach2015lrp}, and DeepLIFT
\citep{shrikumar2017deeplift} are, as published, also defined against architectural assumptions
--- clean gradient flow through a single tower, or a relevance-conservation rule defined layer by
layer --- that a fused-modality encoder like Florence-2's does not straightforwardly satisfy
either, and none of the three was implemented or verified against this architecture within the
pilot's scope. We recommend the paper state plainly that the cross-attention substitution made
here is a documented, verified-correct \emph{choice among constraints}, not evidence that
rollout-style or gradient-based attribution is impossible for fused encoders in general --- only
that none of the textbook implementations applies without the same kind of architectural
verification this chapter performed, and that verification, not assumption, is what should
precede any attribution method's adoption against a new architecture. This is also the natural
place to surface Chapter 14's open question~1 to Professor Abdallah, on which attribution method,
if any, the paper should treat as acceptable evidence for a fused-modality vision-language model
when none of the standard methods applies in its textbook form.

Separately, the size confound identified in this chapter's central finding should be carried
forward as a constraint on any cross-rule sparsity comparison later in the pilot:
\textbf{comparing sparsity across rules without controlling for the grounded object's physical
size will make large-object rules (rule~4's excavator) look artificially worse than small-object
rules (rule~1's hard hat) by this metric alone}, independent of any real difference in explanation
quality. Finally, the design decision to report visual and text attribution streams
\textbf{separately, never merged into one score}, is carried forward from this chapter rather
than introduced fresh later --- the same separation is mirrored by Chapter~9's answer-level versus
explanation-level robustness split, and recording the precedent here lets that later split read
as a continuation of an established practice rather than as a new and unmotivated one.
